\chaves{Escreva aqui as palavras-chave, separadas por vírgula. Exemplo: 
engenharia de software; experiência prática; desenvolvimento de software; manutenção de software.
}

\begin{resumo} 
O resumo estruturado do trabalho deve ter uma página, no formato do seguinte exemplo:

\textbf{Contexto:}
Implementação e testes de software são atividades essenciais de engenharia de software, garantindo que funções sejam corretamente desenvolvidas e que o software atenda aos requisitos definidos. Essas atividades exigem integração entre ferramentas, práticas de codificação e estratégias de teste que assegurem qualidade e confiabilidade.

\textbf{Objetivo:}
Este trabalho tem como objetivo relatar a experiência adquirida durante a implementação de código e a execução de testes em um projeto de software, destacando técnicas aplicadas, desafios enfrentados e contribuições para melhoria de processos.

\textbf{Método:}
A experiência foi baseada na metodologia de pesquisa-ação, no contexto de um projeto conduzido em uma fábrica de software, utilizando desenvolvimento incremental e ferramentas de automação de testes. O relato baseou-se em análise de registros de atividades, documentação do projeto e observações realizadas ao longo do processo.

\textbf{Resultados:}
A experiência resultou na implementação bem-sucedida do projeto, com detecção e correção de falhas por meio de testes sistemáticos. As ferramentas de automação contribuíram para a rastreabilidade dos testes. 
A experiência permitiu consolidar e aprimorar habilidades em implementação e teste de software. 

\textbf{Conclusões:}
Os resultados demonstram que a integração entre implementação e testes promove a qualidade do software. O trabalho permitiu desenvolver habilidades práticas em testes, compreender desafios de implementação e reconhecer boas práticas de engenharia de software. A experiência sugere que ampliar testes automatizados e o uso de métricas de qualidade são uma oportunidade de melhoria para novas iterações do projeto.

\end{resumo}
